\section{Atomic Operations}

\subsection{Why atomics exist}

Consider this code running with interrupts or multiple execution contexts:

\begin{lstlisting}[language={}]
counter = counter + 1;
\end{lstlisting}

At machine level, that is usually:

\begin{lstlisting}[language={}]
lw   t0, 0(a0)
addi t0, t0, 1
sw   t0, 0(a0)
\end{lstlisting}

If an interrupt or another hart modifies \texttt{counter} between the load and store, an update can be lost. Atomic instructions solve this by making read-modify-write sequences indivisible at the architectural level.

\subsection{RISC-V atomic concepts}

RISC-V atomics include two styles:

\begin{enumerate}
\item \textbf{AMO instructions}: read old memory value, compute a new value, write it back, and return the old value.
\item \textbf{LR/SC instructions}: load-reserved and store-conditional. \texttt{lr.w} creates a reservation. \texttt{sc.w} stores only if the reservation is still valid.
\end{enumerate}

\subsection{Implementation in this repository}

\texttt{cpu\_core.v} contains atomic datapath logic for opcode \texttt{0101111}:

\begin{itemize}
\item \texttt{funct7[6:2]} is used as the AMO function field;
\item \texttt{lr.w} records a reservation address;
\item \texttt{sc.w} succeeds only if the reservation is valid and matches the address;
\item AMO read-modify-write instructions compute \texttt{amo\_alu} from memory data and \texttt{rs2};
\item atomics are word-sized in this design (\texttt{dmem\_funct3 = 3'b010}).
\end{itemize}
The reservation state is held in:

\begin{lstlisting}[language={}]
resv_valid
resv_addr
\end{lstlisting}

The design clears the reservation after store-conditional and AMO behavior, and also clears it on traps in the relevant logic. That matters because a task switch between \texttt{lr.w} and \texttt{sc.w} should normally cause the later \texttt{sc.w} to fail.

\subsection{Important limitation}

The atomic path is present in the SoC-style core logic, not uniformly integrated into every pipeline and MMU variant. A pipelined implementation would need to ensure the read-modify-write sequence is not split incorrectly across stages or interrupted by hazards.

\bigskip
